% Part: first-order-logic
% Chapter: model-theory
% Section: theory-of-m

\documentclass[../../../include/open-logic-section]{subfiles}

\begin{document}

\olfileid{mod}{bas}{thm}

\section{কোনো \printtoken{S}{structure}-এর তত্ত্ব}

প্রত্যেক !!{structure}~$\Struct{M}$ কিছু !!{sentence}s-কে সত্য এবং
কিছু বাক্যকে মিথ্যা করে। সেটি যেসব !!{sentence}s-কে সত্য করে, তাদের
সকলের সেটকে তার \emph{তত্ত্ব} বলা হয়। সেটটি সত্যিই একটি তত্ত্ব, কারণ
সেটি থেকে অনুসৃত সব কিছুই তার সব মডেলে সত্য হতে হবে, বিশেষত
$\Struct{M}$-এও।

\begin{defn}
  কোনো !!a{structure}~$\Struct M$ দেওয়া থাকলে, $\Struct{M}$-এর
  \emph{তত্ত্ব} হলো $\Struct{M}$-এ সত্য !!{sentence}s-এর সেট
  $\Theory{M}$; অর্থাৎ $\Theory{M} = \Setabs{!A}{\Sat{M}{!A}}$।
\end{defn}

অনানুষ্ঠানিকভাবে আমরা “তত্ত্ব” শব্দটি অভিপ্রেত ব্যাখ্যাসহ
!!{sentence}s-এর সেট বোঝাতেও ব্যবহার করি, সেটটি নিগমনগতভাবে বদ্ধ হোক
বা না হোক।

\begin{prop}
যেকোনো $\Struct{M}$-এর জন্য $\Theory{M}$ সম্পূর্ণ।
\end{prop}

\begin{proof}
যেকোনো !!{sentence}~$!A$-এর জন্য হয় $\Sat{M}{!A}$, নয়
$\Sat{M}{\lnot !A}$; তাই হয় $!A \in \Theory{M}$, নয়
$\lnot !A \in \Theory{M}$।
\end{proof}

\begin{prop}\ollabel{prop:equiv}
  প্রত্যেক $!A \in \Theory{M}$-এর জন্য $\Struct{N} \models !A$ হলে
  $\Struct{M} \elemequiv \Struct{N}$।
\end{prop}

\begin{proof}
সব $!A \in \Theory{M}$-এর জন্য $\Sat{N}{!A}$ বলে $\Theory{M}
\subseteq \Theory{N}$। যদি $\Sat{N}{!A}$ হয়, তবে
$\Sat/{N}{\lnot !A}$; তাই $\lnot !A \notin \Theory{M}$। যেহেতু
$\Theory{M}$ সম্পূর্ণ, তাই $!A \in \Theory{M}$। অতএব $\Theory{N}
\subseteq \Theory{M}$; ফলে $\Struct{M} \elemequiv \Struct{N}$।
\end{proof}

\begin{rem}\ollabel{remark:R}
  $\Struct{R} = \langle\Real, <\rangle$ বিবেচনা করো। এটি এমন একটি
  !!{structure}, যার সংজ্ঞাক্ষেত্র বাস্তব সংখ্যার সেট~$\Real$ এবং যার
  !!{language}-এ কেবল একটি দ্বি-স্থানীয় !!{predicate} আছে, সেটি বাস্তব
  সংখ্যার উপর~$<$ সম্পর্ক হিসেবে ব্যাখ্যাত। স্পষ্টতই $\Struct{R}$
  !!{nonenumerable}; কিন্তু $\Theory{R}$ স্পষ্টতই সংগত বলে
  লোয়েনহাইম--স্কোলেম উপপাদ্য অনুসারে তার একটি !!a{enumerable} মডেল
  আছে, ধরা যাক $\Struct{S}$। \olref{prop:equiv} অনুসারে
  $\Struct{R} \equiv \Struct{S}$। আরও, $\Struct{R}$ ও~$\Struct{S}$
  সমরূপ নয় বলে এটি দেখায় যে \olref[iso]{thm:isom}-এর বিপরীতটি
  সাধারণভাবে সত্য নয়।
\end{rem}

\end{document}
